\documentclass{article}

% page layout
\usepackage[letterpaper, textwidth=125mm, textheight=230mm]{geometry}
\usepackage{setspace}
\setstretch{1.08}
\sloppy\sloppypar\raggedbottom\frenchspacing

% list management
\usepackage{enumitem}
\setlist{itemsep=0.0ex}
\setlist[enumerate]{label=\textsl{(\alph*)}}

% problem typesetting -- THIS IS SO INSANE
\usepackage{amsmath, amsthm}
\makeatletter
\def\th@definition{%
  \thm@notefont{}% same as heading font
  \normalfont % body font
}
\makeatother
\usepackage{etoolbox}
\makeatletter
\patchcmd{\@thm}{\thm@headpunct{.}}{\thm@headpunct{:}}{}{}
\makeatother
\theoremstyle{definition}
\newcommand{\problemname}{Problem}
\newtheorem{problem}{\problemname}
\newcommand{\probref}[1]{\problemname~\ref{#1}}
\theoremstyle{remark}
\newcommand{\notename}{Note}
\newtheorem{note}{\notename}[problem]
\newcommand{\noteref}[1]{\textit{\notename}~\ref{#1}}

% math macros
\newcommand{\dd}{\mathrm{d}}

\begin{document}
\section*{Crowded-field and confusion problem set}
\noindent
\textit{by} \textbf{David W. Hogg}

\medskip\noindent
There are some disagreements between Hogg and Casey about what is and isn't possible in modeling crowded fields.
I (Hogg) am talking out of my ass; this problem set is designed to make my conjectures explicit, and to test them.

\begin{problem}[power-law flux distributions]
In terms of \emph{apparent flux} $S$, sources tend to be drawn from power-laws of the form
\begin{align}
    \frac{\dd\ln n}{\dd\ln S} = A - \alpha\,\ln\frac{S}{S_0} ~,
\end{align}
where $A$ is a log amplitude, $\alpha$ is a power law exponent, and $S_0$ is a characteristic flux.
Write code to draw flux values $S$ from the implied probability distribution, and demonstrate by making plots that your code is correct.
Your code should take as input the power-law exponent $\alpha$, a flux interval $[S_{\min}, S_{\max}]$ flux in which to work, to and the number $n$ of stars you want returned in that interval.

Now what exponent $\alpha$ do you expect for bright stars in the nearby Galaxy?
Make the assumption that stars uniformly fill three-dimensional space and that the luminosity distance and volume elements are trivially Euclidean.

What exponent $\alpha$ do you expect for the stars in a globular cluster, at the faint end, roughly?

How do your answers transform if you ask about the log number per \emph{magnitude} instead of log flux?

Finally (and you will need this for \probref{prob:fullposterior}), what are the corresponding probability densities $p(S)$ and $p(\ln S)$ corresponding to this power law?
The normalizations of these pdfs will depend on all three control parameters $\alpha, S_{\min}, S_{\max}$.
\end{problem}

\begin{problem}[naive detectability and uncertainties for an isolated star]\label{prob:naive}
Imagine that you have a background-limited image taken with an ideal instrument consisting of a square grid of identical pixels.
Background-limited, here, means that every pixel has a noise contribution that doesn't depend significantly on the sources that contribute to that pixel's intensity.
Imagine that the pixel-convolved point-spread function (PRF) in this image is precisely Gaussian and has a full width at half maximum (FWHM) $w=3.0$\,pixels.
Let's imagine that we took an image that contains, within its field of view, exactly one unresolved star of flux $S$, and the star is nowhere near the edge of the imaging field of view.
Finally, assume that there is a white (uncorrelated, iid) noise contribution to every pixel with value $\sigma=1.$ (in whatever units).
Assume that there is perfect and uniform detector sensitivity.
At what flux $S_5$ is this star measured at $5$-sigma?
Equivalently, what is the expected uncertainty $\sigma_S$ on the brightness of this star?

Give your answers in terms of the noise level $\sigma$ and the PRF FWHM $w$.
That is, I gave values for $\sigma$ and $w$ but give your answers symbolically.
Get as precise or as approximate as you like; discuss the subtle issues.

Now what is the expected uncertainty $\sigma_\theta$ on the \emph{position} (within the image, in pixels) of the star?
Give your answer in terms of $S$, $\sigma$, and $w$.
\end{problem}

\begin{note}
The pixel-convolved point-spread function, or, in astronomer parlance, PRF, is the thing that can be purely \emph{sampled} (that is, not \emph{integrated}) at the pixel-center locations.
\end{note}

\begin{problem}[Generate crowded fields]\label{prob:simulate}
Make a $100\times 100$ pixel grid and fill it with stars.
Here are the requirements:
Generate stellar fluxes $S$ with the power law $\alpha=1$ (the barely diverging, barely Olbers-paradox case).
Generate stellar positions uniformly in the image, but generate in a larger region, so the stars extend naturally beyond the edges of the image.
Generate stars in the flux interval $[S_{\min},S_{\max}]=[0.003\,S_5, 3.0\,S_5]$ where $S_5$ is the flux level you computed in \probref{prob:naive}.
That is, generate stars down to extremely low fluxes such that the faintest would be ``detected'' at the 0.015-sigma level (of course they won't be detected).
Generate enough stars that there are of order 300 stars at true fluxes $S>S_5$ within the image.
Make a synthetic image using the PRF stated in \probref{prob:naive} and add uniform white noise as required by \probref{prob:naive}.
Keep (don't delete) the true catalog that went into the image!
Deliver this image and this catalog, and code that can generate arbitrary numbers of new examples.
\end{problem}

\begin{note}
Why $300$? Because that corresponds to about 30 beams per source.
\end{note}

\begin{note}
You will have to generate a \emph{fuck ton} of stars to meet the requirements of this problem.
\end{note}

\begin{problem}[crowding contributions to flux and position errors]\label{prob:trad}
Use whatever standard astrophysical code you would use to detect stars in the generated image.
Use the best (frequentist, standard) crowded-field code to which you have access; don't skimp here.
Ask the code to deliver all stars down to the flux level $S_5$ you computed in \probref{prob:naive}.
Document your choices.
Compare this catalog (the catalog generated by the detection code) to the (bright end of the) true catalog that went into the image generation.
How do the differences compare to the uncertainty expectations you generated in \probref{prob:naive}?

Run again but generating an image that only has one $S>S_5$ source in the field.
That is, make a less crowded case.
Are your two catalogs more similar in this case?
Comment.
For extra credit, plot the variances of brightness and position of $S\sim S_5$ sources as a function of the density of $S>S_5$ sources (in terms of sources per beam, perhaps).
\end{problem}

\begin{note}
If you do the extra-credit part, you will have to define the number of ``beams'' in the image, given the PSF FWHM $w$.
You have choices here; make them sensibly.
\end{note}

\begin{problem}[likelihood function]\label{prob:lf}
Make a likelihood function for a catalog that spans your image.
That is, write a function that takes in a noisy image like you made in \probref{prob:simulate}, a catalog of stars with positions $(x,y)_j$ (in pixel coordinates) and fluxes $S_j$, and a sky level $I_0$, and outputs a log likelhood value.
This log likelihood should be the log of a probability density for the image data, given the catalog and sky level.
When you construct this likelihood function, use the \emph{true, non-wrong} values of the per-pixel white noise $sigma$ and PRF FWHM $w$.
That is, make the likelihood function \emph{exactly correct}, which is an unrealistically optimistic setting for real data.

Figure out how to \emph{test} your likelihood function code.
That is, how do you convince an external customer that you are computing the correct likelihood?
One option is to ``slice'' the likelihood, or plot the likelihood as you vary one parameter at a time, for one of your bright sources?
What expectations do you have for such dependencies?
Your answers to \probref{prob:naive} should be very useful here.

Very importantly, your likelihood function must include a parameter for the ``dc'' or ``piston'' or sky or background level $I_0$ of the image.
You are not permitted to assume that you have data for which this has been calibrated to zero (that's not possible, ever).
\end{problem}

\begin{note}
Be sure to permit the catalog to include sources that are outside the field of view of the image. After all, there are always stars at, and outside, the image boundary, and the PRF has broad support.
\end{note}

\begin{note}
Once you have this likelihood function and it is executable, then you have a forward model for your data and you can perform principled inferences.
Let's go!
\end{note}

\begin{problem}[maximum-likelihood catalog]\label{prob:mlcatalog}
Choose one of the images you generated in \probref{prob:simulate} as your ``data.''
By being extremely and outrageously clever, figure out how to make the maximum-likelihood catalog and maximum-likelihood sky level for this image.
That is, optimize the full likelihood function.
Note how hard this is: The catalog can have variable numbers of sources, and the likelihood function contains a combinatorial number of exact degeneracies (think: permutations), and boatloads of near-degeneracies.
If you achieve anything here, draft a paper about it for the refereed literature.

Compare your results to the traditional results you got in \probref{prob:trad}.
Is your maximum-likelihood method any better than standard practice?

You will be considered to have solved this problem even if you only solve it for tiny image patches centered on the $S>S_5$ stars.
That is, you don't have to successfully optimize everything to be deemed successful here.
But you \emph{do} need to consider catalog sources outside the target image patches.
\end{problem}

\begin{note}
If I were forced to do this problem, I would start with something that looks like annealing, possibly with split-and-merge moves.
Or else I would put down a grid of very faint sources at fixed locations, and only vary the \emph{fluxes} of the faint sources.
Only the bright sources would be given full positional and flux freedom.
Would that work?
\end{note}

\begin{note}
I conjecture that maximum-likelihood will not give better answers than standard practice.
After all, standard practice is an approximation to maximum-likelihood.
\end{note}

\begin{problem}[fully profiled catalog]
Once you have solved \probref{prob:mlcatalog} it would be possible to get \emph{profile} likelihoods for the $S>S_5$ sources, profiling over the fainter (or, really \emph{all} other) sources.
That would be incredible, and should be part of the paper you draft as part of the solution to \probref{prob:mlcatalog}.
\end{problem}

\begin{note}
The fully profiled catalog might be just as valuable as the fully marginalized catalog you will obtain below in \probref{prob:fullposterior}.
\end{note}

\begin{problem}[sources outside the image]
A side quest is to show that your model or your likelihood optimization \emph{requires}, statistically, sources outside the image footprint.
If you succeed, draft an April Fools' paper about this subject.
\end{problem}

\begin{problem}[full marginalization in catalog space]\label{prob:fullposterior}
Write log-prior-pdf functions that describe our expectations for the positions and fluxes of the catalog sources.
For positions, use uniforms (that extend sufficiently outside the image to deal with edge stars).
For fluxes, use the \emph{exactly correct} power-law distribution with \emph{exactly correct} minimum and maximum fluxes that you used in \probref{prob:simulate} to create the data in the first place.
For the sky level $I_0$ choose something sensible (or, for extra credit, something \emph{correct}, given the generative model).
These log-prior functions should add to give you an executable prior pdf for any catalog plus sky level, independent of the data.

Now combine this log-prior function with your log-likelihood function from \probref{prob:lf} to create an incorrectly normalized (but sufficient for a log-posterior function input to an MCMC method).
This log-posterior function takes in a catalog and sky level, and returns either (if the log prior is negative infinity) negative infinity or else a finite value that is offset by a constant from the true log posterior probability density.

Now figure out how to sample the shit out of this with some kind of MCMC method to develop fully marginalized posterior beliefs about the catalog.
Are the resulting posterior beliefs for individual sources more accurate than what you got with traditional methods in \probref{prob:trad} or with your maximum-likelihood method in \probref{prob:mlcatalog}?

You will be considered to have solved this problem if you can solve it for tiny image patches centered on only the $S>S_5$ stars.
That is, you don't have to successfully solve all of Bayesian inference simultaneously over the full image footprint to be seen as solving this problem.
Even if you solve this very limited version of the problem, draft a paper for the refereed literature about it.
\end{problem}

\begin{note}
I conjecture that the result of the Bayesian marginalization will not be better than the traditional results found in \probref{prob:trad}.
If they are, this is very big news.
The results are publishable either way.
\end{note}

\begin{note}
There are approaches here with reversible-jump MCMC.
There are also approaches with ``massed priors.''
If I were forced to take the shot, I would use the latter.
\end{note}

\begin{note}
I conjecture that this problem is close to impossible to solve at finite autocorrelation time with any known sampler.
\end{note}

\end{document}
